Penelitian ini bertujuan untuk membangun dataset deteksi objek sembilan 
kelas khas kawasan Malioboro, membandingkan dan mengoptimasi varian model 
YOLO untuk mendeteksi kelas-kelas tersebut, lalu mengimplementasikan model 
terpilih ke dalam sistem pemantauan kepadatan pengunjung berbasis web, 
sebagai jawaban atas keterbatasan pemantauan manual dan ketidaktersediaan 
dataset publik seperti COCO untuk objek lokal seperti andong, becak, dan 
bajaj. Metodologi yang digunakan meliputi pembangunan dataset dari rekaman 
22 kamera CCTV publik kawasan Malioboro melalui anotasi semi-otomatis 
berbantuan model \textit{pretrained} dengan koreksi manual, pelatihan dan 
perbandingan delapan varian model YOLOv8 dan YOLO11 menggunakan metrik 
mAP50, mAP50-95, dan F1-\textit{score} pada himpunan uji, penyetelan 
resolusi citra masukan pada model terbaik, serta implementasi model 
tersebut ke dalam sistem berbasis web yang memproses \textit{stream} video 
kamera CCTV secara \textit{real-time} dengan ambang notifikasi kepadatan 
berdasarkan konsep tingkat pelayanan (\textit{level of service}) pejalan 
kaki. Hasil penelitian menunjukkan dataset akhir memuat 1.581 citra 
beranotasi dengan distribusi kelas yang sangat tidak seimbang, mencapai 
rasio 237 berbanding 1 antara kelas terbanyak dan tersedikit. Pada himpunan 
uji, yolov8l mencatat mAP50 dan rerata F1 tertinggi (0,6175 dan 0,6222), 
sedangkan yolo11l mencatat mAP50-95 tertinggi (0,4012) dengan ukuran model 
42\% lebih ringkas sehingga dipilih sebagai model terbaik. Penyetelan 
resolusi dari 640 menjadi 960 piksel kemudian menaikkan rerata F1 model 
terpilih dari 0,5962 menjadi 0,6739, melampaui seluruh varian pada resolusi 
dasar. Pengujian pada perangkat penerapan menunjukkan model terpilih 
mencapai 9,73 \textit{frame} per detik pada GPU GTX 1650, lebih cepat 
dibanding yolov8l yang hanya 8,80 \textit{frame} per detik, dengan sistem 
yang berhasil menampilkan deteksi langsung, grafik kepadatan, dan riwayat 
peringatan sesuai rancangan. Kesimpulan dari penelitian ini adalah model 
YOLO11 \textit{large} dengan resolusi masukan yang disetel terbukti menjadi 
pilihan yang seimbang antara akurasi, ukuran, dan kecepatan untuk 
diterapkan pada sistem pemantauan kepadatan berbasis web di kawasan 
Malioboro, dengan ketidakseimbangan kelas pada dataset lokal tetap menjadi 
faktor yang memengaruhi performa deteksi pada kelas minoritas.

\noindent\textbf{Kata kunci} : deteksi objek, YOLO, kepadatan pengunjung, ketidakseimbangan kelas, Malioboro